\section{Discussion}
\label{sec:discussion}

Dynamic agent regimes provide a precise way to ask how changing model selection changes the outcome of an agent's complete execution. The central requirement is to align the intervention, the trajectory distribution, and the evaluation design. Versioned model configurations define the available actions; prospective eligibility defines when those actions can change; and sequential assignment probabilities connect the collected trajectories to supported target policies. Under the stated assumptions, the analysis reduces to longitudinal policy evaluation. The additional results clarify how that evaluation relates to selectively observed branch continuations and to bounded changes in the execution mechanism.

This formulation does not replace existing dynamic-regime or off-policy methods. The g-formula, importance ratios, fixed-policy influence function, and doubly robust construction have established foundations. Likewise, variance-based allocation and simulation-lemma arguments precede this application. The methodological value of the present development is the explicit connection of these components to decisions that can actually be deployed and outcomes that can actually be observed. Its practical value will depend on whether those contracts can be maintained and whether the resulting estimates are informative at realistic sample sizes.

Overlap remains a central limitation. Randomization can provide support for a prespecified model pool, but a small probability of a target-compatible action sequence can still produce unstable estimates. Restricting routing to fewer eligible opportunities can shorten the likelihood-ratio product, at the cost of restricting the strategies under consideration. Supported stochastic changes similarly narrow the question; they do not reveal the value of an unobserved model or guarantee a useful improvement. Neither an effective-sample-size diagnostic nor a fitted continuation model repairs missing support. High-dimensional histories also make nuisance estimation difficult, and a convenient compressed representation need not preserve the identifying information.

The branch results require their own design contract. The selected prefix must be restored faithfully, selection probabilities must be known, and the prefix weight must correspond to the target history distribution. Pairing continuations can improve precision when it induces favorable covariance, but does not make descendants independent observations. Unequal numbers of branches can change the estimand if their contributions are averaged without appropriate weighting. Sampling structure is consequential in off-policy evaluation more generally \citep{kallus2021multiple}. The derived allocation is optimal only within the stated selection class, for the specified cost model and conditional moments. Estimating these moments in a pilot creates an additional design problem; the oracle result is not a guarantee for an adaptively tuned implementation.

The execution-drift bound has a different role from identification. It propagates assumed uniform kernel discrepancies to a worst-case value discrepancy. It does not show that model updates have small discrepancies, and checking a finite set of restored histories cannot certify a uniform bound over all target-reachable histories. A coarse reward range or accumulated discrepancy may make the bound uninformative. Nevertheless, explicitly stating the required bound distinguishes a sensitivity argument from an unsupported assertion that an earlier estimate remains valid after changing prompts, checkpoints, quantization, tools, or serving behavior.

Several practical conditions also remain outside the current model. The analysis assumes independent task units, with repeated executions and branches handled within task. Shared mutable environments, caches, or resource contention can introduce interference beyond that structure. Outcomes must correspond to the declared verifier and horizon; policy-dependent missing verification requires additional observation assumptions. A fixed benchmark supports conclusions conditional on its tasks unless a task-sampling argument is supplied. Even when policy-value inference is valid, it does not identify which model is best for each individual task or justify unrestricted search over policies using the same evaluation outcomes.

The next empirical step is to compare offline estimates with independently executed, frozen policies under a common contract. Known-truth simulations should separate algebraic correctness from coverage, support failure, and nuisance misspecification. Live studies should use models capable of nondegenerate task performance, record actual routing probabilities, and distinguish task sampling from repeated-seed variability. Restored same-model controls are needed to assess branch reproducibility. Both null improvement and cases in which direct policy execution is cheaper or more precise than off-policy evaluation are scientifically informative. The present theoretical development does not establish empirical superiority or deployment savings; the planned studies are required to determine when this evaluation design is useful.
